% Auto-generated from meeting_2566-09-15_curriculum-draft-1.md (กบ.วช. format v2)
\documentclass{meetingminutes}
\meetingcommittee{คณะกรรมการบริหารหลักสูตร วิศวกรรมคอมพิวเตอร์และปัญญาประดิษฐ์ (บัณฑิตศึกษา)}
\meetingnumber{๑}{๒๕๖๖}
\meetingdate{วันที่ ๑๕ กันยายน พ.ศ. ๒๕๖๖}
\meetingplace{ห้องประชุมคณะเทคโนโลยีอุตสาหกรรม}
\meetingstart{๐๙.๐๐ น.}
\meetingend{๑๑.๔๕ น.}

\begin{document}
\maketitle

\psection{ผู้มาประชุม}
\begin{participants}
  \pmember{1}{รองศาสตราจารย์ ดร.อิสระ อินจันทร์}{ที่ปรึกษาหลักสูตร}{ประธาน}
  \pmember{2}{รองศาสตราจารย์ ดร.กันต์ อินทุวงศ์}{ที่ปรึกษาหลักสูตร}{รองประธาน}
  \pmember{3}{ผู้ช่วยศาสตราจารย์ ดร.พิศิษฐ์ นาคใจ}{อาจารย์ประจำหลักสูตร}{กรรมการ}
  \pmember{4}{ผู้ช่วยศาสตราจารย์ ดร.กาญจนา ดาวเด่น}{อาจารย์ประจำหลักสูตร (เลขานุการ)}{กรรมการ/เลขานุการ}
  \pmember{5}{ผู้ช่วยศาสตราจารย์ ดร.พัชรี มณีรัตน์}{อาจารย์ประจำหลักสูตร}{กรรมการ}
  \pmember{6}{ผู้ช่วยศาสตราจารย์ ดร.พิทักษ์ คล้ายชม}{รองคณบดีฝ่ายวิชาการ คณะเทคโนโลยีอุตสาหกรรม}{กรรมการ}
  \pmember{7}{ผู้ช่วยศาสตราจารย์ ดร.โสกณ วิริยะรัตนกุล}{อาจารย์ประจำหลักสูตร}{กรรมการ}
  \pmember{8}{ผู้ช่วยศาสตราจารย์ ดร.ธัชชัย อยู่ยิ่ง}{อาจารย์ประจำหลักสูตร}{กรรมการ}
  \pmember{9}{ผู้ช่วยศาสตราจารย์ ดร.อภิศักดิ์ พรหมฝ่าย}{อาจารย์ประจำหลักสูตร}{กรรมการ}
\end{participants}

\psection{ผู้ไม่มาประชุม}
\noindent ไม่มี\par

\startmeeting
\openingchair{รองศาสตราจารย์ ดร.อิสระ อินจันทร์}{กล่าวเปิดการประชุมและดำเนินการประชุมตามระเบียบวาระ ดังนี้}

\agenda{๑}{เรื่องแจ้งให้ที่ประชุมทราบ}
ประธานแจ้งที่ประชุมทราบว่า มหาวิทยาลัยราชภัฏอุตรดิตถ์มีนโยบายขยายการเปิดหลักสูตรในระดับบัณฑิตศึกษาให้ครอบคลุมสาขาวิชาที่ตอบสนองต่อความต้องการของตลาดแรงงานในยุคดิจิทัล โดยเฉพาะด้านปัญญาประดิษฐ์ (AI) และวิศวกรรมคอมพิวเตอร์ ซึ่งมีความต้องการบุคลากรสูงในภาคอุตสาหกรรมและภาครัฐ คณะเทคโนโลยีอุตสาหกรรมได้รับมอบหมายให้เป็นเจ้าภาพหลักในการพัฒนาหลักสูตรนี้ ร่วมกับคณะวิทยาศาสตร์และเทคโนโลยี


\agenda{๒}{เรื่องรับรองระเบียบวาระ}
ที่ประชุมรับรองระเบียบวาระการประชุมโดยไม่มีการเปลี่ยนแปลง


\agenda{๓}{เรื่องสืบเนื่อง — ไม่มี}

\agenda{๔}{เรื่องเสนอเพื่อพิจารณา}
\subagenda{๔.๑}{แนวทางการจัดทำหลักสูตรวิศวกรรมศาสตรมหาบัณฑิต}
ผศ.ดร.พิศิษฐ์ นำเสนอแนวคิดเบื้องต้นของหลักสูตร วศ.ม. สาขาวิชาวิศวกรรมคอมพิวเตอร์และปัญญาประดิษฐ์ ดังนี้

\begin{thailist}
  \item ชื่อหลักสูตร: หลักสูตรวิศวกรรมศาสตรมหาบัณฑิต สาขาวิชาวิศวกรรมคอมพิวเตอร์และปัญญาประดิษฐ์ (Master of Engineering Program in Computer Engineering and Artificial Intelligence)
  \item ปริญญา: วิศวกรรมศาสตรมหาบัณฑิต (วิศวกรรมคอมพิวเตอร์และปัญญาประดิษฐ์) / M.Eng. (Computer Engineering and Artificial Intelligence)
  \item จำนวนหน่วยกิต: ไม่น้อยกว่า 36 หน่วยกิต
  \item ระยะเวลาการศึกษา: 2 ปีการศึกษา
  \item โครงสร้างหลักสูตร: 2 แผนการศึกษา คือ แผน 1 (วิชาการ — เน้นวิทยานิพนธ์) และแผน 2 (วิชาชีพ — เน้นสารนิพนธ์)
\end{thailist}
ที่ประชุมอภิปรายและเห็นชอบในหลักการ โดยมีข้อเสนอแนะให้ศึกษาหลักสูตรอ้างอิงจากสถาบันการศึกษาชั้นนำในประเทศ และออกแบบเนื้อหาให้สอดคล้องกับมาตรฐานคุณวุฒิระดับบัณฑิตศึกษาของ กมอ. พ.ศ. 2565

\subagenda{๔.๒}{กรอบผลลัพธ์การเรียนรู้ที่คาดหวัง (PLOs)}
ที่ประชุมอภิปรายกรอบ PLOs เบื้องต้น โดยพิจารณาจาก: - มาตรฐานคุณวุฒิระดับบัณฑิตศึกษา กมอ. 2565 (4 ด้าน: ความรู้/ทักษะ/จริยธรรม/ลักษณะบุคคล) - Bloom's Taxonomy ระดับสูง (Analyze, Evaluate, Create) - แนวโน้มอุตสาหกรรม AI และเทคโนโลยีดิจิทัล

กรอบ PLOs เบื้องต้นที่ที่ประชุมเสนอมี 5 ด้าน ได้แก่ ด้านการสร้างสรรค์นวัตกรรม ด้านการประยุกต์ใช้ความรู้ ด้านการวิจัย/สร้างองค์ความรู้ ด้านจริยธรรมวิชาชีพ และด้านทักษะชีวิตและการเรียนรู้ตลอดชีวิต

\resolution{รับทราบกรอบ PLOs เบื้องต้น และมอบหมายให้คณะทำงานยกร่าง PLO ฉบับสมบูรณ์ พร้อมระบุระดับ Bloom's Taxonomy ให้ชัดเจนในการประชุมครั้งถัดไป}
\subagenda{๔.๓}{กำหนดการสำรวจความต้องการผู้มีส่วนได้ส่วนเสีย}
ที่ประชุมเห็นชอบให้ดำเนินการสำรวจความต้องการผู้มีส่วนได้ส่วนเสีย (Stakeholder Needs Analysis) เพื่อใช้เป็นฐานในการออกแบบ PLOs ตามข้อกำหนด AUN-QA R1.4 โดยกำหนดกลุ่มเป้าหมายการสำรวจ ได้แก่ ผู้ประกอบการภาคอุตสาหกรรม, ภาครัฐและรัฐวิสาหกิจ, สถาบันการศึกษา, และผู้สนใจเข้าศึกษา

\resolution{มอบหมายให้ ผศ.ดร.กาญจนา ดาวเด่น ออกแบบแบบสอบถามสำรวจความต้องการ และดำเนินการเก็บข้อมูลภายในภาคเรียนที่ 2 ปีการศึกษา 2566}

\agenda{๕}{เรื่องอื่น ๆ}
กำหนดการประชุมครั้งต่อไป: เดือนตุลาคม 2566 (วันและเวลาจะแจ้งให้ทราบภายหลัง)


\adjournmeeting

\begin{sigblock}
\typist{ผู้ช่วยศาสตราจารย์ ดร.กาญจนา ดาวเด่น}
\reviewer{รองศาสตราจารย์ ดร.อิสระ อินจันทร์}{กรรมการ}
\chairsig{รองศาสตราจารย์ ดร.อิสระ อินจันทร์}{ประธานการประชุม}
\end{sigblock}

\end{document}
